% =============================================================
% บทที่ 2 - งานวิจัยและทฤษฎีที่เกี่ยวข้อง
% โครงงาน: การออกแบบและพัฒนาเกมเดินเรือด้วยระบบสั่งการด้วยเสียงเพื่อส่งเสริมทักษะการสื่อสาร
% =============================================================

% -----------------------------------------------------------
\section{งานวิจัย โครงงาน และระบบงานที่เกี่ยวข้อง}
% -----------------------------------------------------------

ในส่วนนี้จะนำเสนอทฤษฎี งานวิจัย และเกมที่มีลักษณะคล้ายคลึงกับโครงงานนี้ เพื่อศึกษาแนวคิด จุดเด่น และข้อจำกัด แล้วนำมาเป็นแนวทางในการออกแบบและพัฒนาระบบ

\subsection{การเรียนรู้ภาษาผ่านเกมด้วยเทคโนโลยีการประมวลผลภาษาธรรมชาติ}

งานวิจัยของ Sadigzade\cite{sadigzade2025} นำเสนอการใช้เกมดิจิทัลร่วมกับเทคโนโลยีการรู้จำเสียงและการประมวลผลภาษาธรรมชาติ เพื่อสนับสนุนการเรียนรู้ภาษาอังกฤษ การให้ผู้เรียนได้โต้ตอบด้วยเสียงกับตัวละคร AI ภายในเกมช่วยสร้างสภาพแวดล้อมจำลองที่สมจริง ผลการศึกษาพบว่าช่วยเพิ่มแรงจูงใจ ลดความวิตกกังวลในการพูด และพัฒนาความคล่องแคล่วในการสื่อสารได้อย่างมีประสิทธิภาพ อย่างไรก็ตาม เกมในลักษณะดังกล่าวมักไม่ได้เน้นความกดดันด้านเวลาหรือสถานการณ์จำลองที่ต้องใช้ไหวพริบอย่างรวดเร็ว

\subsection{เครื่องมือเกมควบคุมด้วยเสียงเพื่อเสริมสร้างการเรียนรู้}

งานวิจัยของ Varma และคณะ\cite{varma2025} มุ่งเน้นพัฒนาเกมที่ควบคุมการทำงานด้วยคำสั่งเสียงแบบเรียลไทม์ โดยทำงานผ่านการแปลงเสียงพูดเป็นคำสั่งเพื่อควบคุมการเคลื่อนที่ขององค์ประกอบในเกมแทนการใช้อุปกรณ์ควบคุมปกติ การควบคุมผ่านเสียงช่วยกระตุ้นให้ผู้เล่นต้องคิดตัดสินใจอย่างรวดเร็วและเปล่งเสียงให้ชัดเจน ซึ่งช่วยพัฒนาทักษะการตอบสนองได้ดี

\subsection{การสำรวจการโต้ตอบด้วยเสียงกับตัวละคร NPC ในเกม}

Zargham และคณะ\cite{zargham2024} ได้สำรวจความคิดเห็นของผู้เชี่ยวชาญเกี่ยวกับการใช้เสียงพูดเพื่อสื่อสารกับตัวละคร NPC ในวิดีโอเกม ผลการวิจัยระบุว่าการสั่งการด้วยเสียงช่วยเพิ่มความรู้สึกดื่มด่ำ (Immersion) และความสมจริง นอกจากนี้ การให้ NPC ตอบสนองต่อคำสั่งเสียงอย่างเป็นธรรมชาติยังช่วยเพิ่มความผูกพันได้ แต่จำเป็นต้องคำนึงถึงบริบทและความแม่นยำของระบบรู้จำเสียงเพื่อป้องกันความหงุดหงิดจากการแปลความหมายผิดพลาด

\subsection{เกมจำลองการควบคุมยานพาหนะด้วยเสียง (Commercial Games)}

ในอุตสาหกรรมเกมเชิงพาณิชย์ ได้มีการนำระบบรู้จำเสียงมาใช้เพื่อเพิ่มความสมจริงในเกมแนวจำลองสถานการณ์ (Simulation) เช่น เกม Powerboat VR\cite{powerboat2020} ที่ผู้เล่นสามารถใช้คำสั่งเสียงภาษาอังกฤษในการควบคุมทิศทางและความเร็วของเรือแทนคอนโทรลเลอร์ปกติ หรือเกม Tower! 3D\cite{tower3d2017} ที่ผู้เล่นต้องใช้เสียงพูดตามมาตรฐานคำศัพท์การบินเพื่อสั่งการเครื่องบิน เกมเหล่านี้แสดงให้เห็นว่าคำสั่งเสียงสามารถสร้างประสบการณ์ที่ดื่มด่ำได้เป็นอย่างดี

\subsection{สรุปเปรียบเทียบงานวิจัยและเกมที่เกี่ยวข้อง}

จากการทบทวนระบบและงานวิจัยข้างต้น สามารถสรุปเปรียบเทียบคุณสมบัติหลักได้ดังตารางที่ \ref{tab:compare_research}

\begin{table}[H]
\centering
\caption{เปรียบเทียบคุณสมบัติของระบบและเกมที่เกี่ยวข้อง}
\label{tab:compare_research}
\begin{tabular}{|p{3.5cm}|p{2cm}|p{2cm}|p{2cm}|p{2cm}|}
\hline
\textbf{ระบบ/เกม} & \textbf{ควบคุมด้วยเสียง} & \textbf{NPC โต้ตอบด้วยเสียง} & \textbf{จำลองฟิสิกส์} & \textbf{ลดความกังวลภาษา} \\
\hline
งานวิจัย NLP \& Games & \cmark & \cmark & -- & \cmark \\ \hline
Voice-Controlled Tools & \cmark & -- & \cmark & -- \\ \hline
Powerboat VR & \cmark & -- & \cmark & -- \\ \hline
In Verbis Virtus & \cmark & -- & -- & -- \\ \hline
\textbf{โครงงานนี้} & \cmark & \cmark & \cmark & \cmark \\
\hline
\end{tabular}
\end{table}

จากตารางที่ \ref{tab:compare_research} พบว่าส่วนใหญ่ระบบจะแยกกันระหว่างเกมที่เน้นการเรียนรู้ภาษาแต่ขาดความสมจริงทางฟิสิกส์ กับเกมจำลองสถานการณ์ที่มีความสมจริงแต่ไม่ได้เน้นกระบวนการให้ข้อมูลตอบกลับ (Feedback) เพื่อแก้ปัญหาความกังวลในการใช้ภาษา โครงงานนี้จึงบูรณาการคุณสมบัติทั้ง 4 ด้านเข้าด้วยกัน

\subsection{เปรียบเทียบเทคโนโลยีเอนจินพัฒนาเกม (Game Engine)}

{\sloppy
ในการเลือกเทคโนโลยีสำหรับพัฒนาเกม ได้เปรียบเทียบ Game Engine ที่นิยม 3 ตัวเลือก ดังตารางที่ \ref{tab:compare_engine}
}

\begin{table}[H]
\centering
\caption{เปรียบเทียบเทคโนโลยี Game Engine}
\label{tab:compare_engine}
\begin{tabular}{|p{3cm}|p{3.5cm}|p{3.5cm}|p{3cm}|}
\hline
\textbf{เกณฑ์} & \textbf{Unity 3D} & \textbf{Unreal Engine} & \textbf{Godot} \\
\hline
ภาษาหลักที่ใช้ & C\# & C++ / Blueprint & GDScript / C\# \\ \hline
การรองรับ 3D & ดีมาก & ยอดเยี่ยม & ดี \\ \hline
การเชื่อมต่อ API ภายนอก & ง่าย มีไลบรารีรองรับเยอะ & ปานกลาง & ปานกลาง \\ \hline
ชุมชนผู้พัฒนา & ใหญ่มาก & ใหญ่มาก & กำลังเติบโต \\ \hline
\end{tabular}
\end{table}

จากตารางที่ \ref{tab:compare_engine} โครงงานนี้เลือกใช้ \textbf{Unity 3D} ร่วมกับภาษา \textbf{C\#} เนื่องจากมีความยืดหยุ่นสูง มีไลบรารีสำหรับการจัดการเสียงที่ครอบคลุม และผู้พัฒนามีความคุ้นเคยในการเชื่อมต่อกับ API ภายนอก

% -----------------------------------------------------------
\section{หลักการและทฤษฎีที่เกี่ยวข้อง}
% -----------------------------------------------------------

\subsection{การเรียนรู้โดยใช้เกมเป็นฐาน (Digital Game-Based Learning: DGBL)}

เป็นการนำเกมดิจิทัลมาใช้เป็นเครื่องมือเพื่อส่งเสริมกระบวนการเรียนรู้ โดยอาศัยความสนุกสนานและความท้าทายเป็นตัวกระตุ้นแรงจูงใจภายใน (Intrinsic Motivation) ทฤษฎีนี้เน้นให้ผู้เรียนเกิดการเรียนรู้เชิงประสบการณ์ (Experiential Learning) โครงงานนี้นำแนวคิดนี้มาเป็นแกนหลักในการเปลี่ยนรูปแบบการฝึกพูดภาษาอังกฤษให้กลายเป็นการเล่นเกมเพื่อบรรลุภารกิจ

\subsection{ทฤษฎีการลดตัวกรองทางอารมณ์ (Affective Filter Hypothesis)}

ทฤษฎีนี้อธิบายว่าตัวกรองทางอารมณ์ เช่น ความกลัว ความวิตกกังวล และความขาดความมั่นใจ เป็นอุปสรรคสำคัญที่ขัดขวางการเรียนรู้และการพูดภาษาที่สอง\cite{sadigzade2025} โครงงานนี้นำทฤษฎีนี้มาประยุกต์ใช้โดยการใช้ตัวละครลูกเรือปัญญาประดิษฐ์ (AI Crew) เป็นผู้รับฟังคำสั่งแทนมนุษย์จริง ซึ่งช่วยลดตัวกรองทางอารมณ์และลดความวิตกกังวลในการใช้ภาษา (Foreign Language Anxiety) ทำให้ผู้เล่นผ่อนคลายและกล้าเปล่งเสียงมากขึ้น

\subsection{ทฤษฎีการเรียนรู้ตามสถานการณ์ (Situated Learning Theory)}

ทฤษฎีนี้เสนอว่าการเรียนรู้จะมีประสิทธิภาพสูงสุดเมื่อเกิดขึ้นในบริบทหรือสถานการณ์ที่ความรู้นั้นถูกนำไปใช้จริง ผู้เล่นจะเรียนรู้ผ่านบริบทการสั่งการทิศทาง การปรับความเร็ว และการแก้ปัญหาเฉพาะหน้าทางทะเล (Maritime Communication) ซึ่งช่วยให้เกิดการจดจำและประยุกต์ใช้ได้เร็วกว่าการท่องจำ

\subsection{เทคโนโลยีการรู้จำเสียง (Speech Recognition)}

การรู้จำเสียงคือกระบวนการแปลงคำสั่งเสียงพูดเป็นข้อความ (Speech-to-Text) โครงงานนี้เลือกประยุกต์ใช้เทคโนโลยีปัญญาประดิษฐ์ที่มีความแม่นยำสูง เช่น Whisper AI เพื่อประมวลผลเสียงภาษาอังกฤษของผู้เล่นแบบเรียลไทม์ และส่งเป็นข้อมูลนำเข้า (Input) สู่ระบบวิเคราะห์คำสั่ง

\subsection{ปัญญาประดิษฐ์และการประมวลผลภาษาธรรมชาติ (OpenAI API)}

โครงงานนี้มีการใช้บริการเชื่อมต่อระบบประมวลผล (API) เช่น OpenAI GPT-4o mini เพื่อใช้ในระบบโต้ตอบของลูกเรือ ระบบปัญญาประดิษฐ์นี้จะรับข้อความที่แปลงจากเสียงพูด มาทำการวิเคราะห์บริบทและความถูกต้อง หากคำสั่งไม่ชัดเจน AI จะจำลองบทสนทนาโต้ตอบกลับเพื่อทวนคำสั่ง
